采用 Beckner、Carbery、Semmes 和 Soria \cite{r-7}, Lemma~1 中对不交分解作符号平均的方法，并使用他们在 Lemma~2 (p.~395) 中的球冠估计。相关球冠分解及平方函数论证见 Bourgain \cite{r-13}, Section~6, (6.32)--(6.41), pp.~179--181。固定充分小的 \(c>0\)，在球面上取互不相交的球冠 \(C_j=\{|\omega-\omega_j|<c\delta\}\)，令 \(h_j(\omega)=e^{-i\delta^{-2}x_j\cdot\omega}\mathbf1_{C_j}(\omega)\)。在 \(T_j\) 上写作 \(Y-x_j=t\omega_j+v\)，其中 \(|t|\le1/2\)、\(|v|<2\delta\)。因为 \(\omega_j\cdot(\omega-\omega_j)=-|\omega-\omega_j|^2/2\)，\(Eh_j(\delta^{-2}Y)\) 中的相位与其在 \(\omega_j\) 处的取值相差至多 \(c^2/4+2c\)。故在那里 \(|Eh_j(\delta^{-2}Y)|\gtrsim\delta^d\)。

按半开单位格点立方体 \(Q\) 将中心分组，其中 \(J_Q=\{j:x_j\in Q\}\)。一组中的所有管都在 \(5Q\) 内，将它伸缩 \(\delta^{-2}\) 倍后包含在半径为 \(C_n\delta^{-2}\) 的一个球内。取独立的对称符号 \(\varepsilon_j\)，令 \(g_\varepsilon=\sum_{j\in J_Q}\varepsilon_j a_j^{1/2}h_j\)。由不相交性，其 \(L^p\) 范数的 \(p\) 次幂由 \(C\delta^d\sum_{j\in J_Q}a_j^q\) 控制。Jensen 不等式给出
\[
 \mathbb E|Eg_\varepsilon|^p\ge
 \left(\sum_{j\in J_Q}a_j|Eh_j|^2\right)^q .
\]
在 \(\delta^{-2}(5Q)\) 上积分，在包含它的球上使用假设中的估计，再换元至 \(Y\)。这给出
\[
 \delta^{dp-2n}\int_{5Q}\left(\sum_{j\in J_Q}a_j\mathbf1_{T_j}\right)^q
 \lesssim L(1/\delta)^{Ap}\delta^d\sum_{j\in J_Q}a_j^q .
\]
其中 \(dp=2n\)。\(5Q\) 的重叠次数有统一上界，故所求范数的 \(q\) 次幂不超过各组积分之和的常数倍。取 \(q\) 次方根并用 \(Ap/q=2A\)。
